\documentclass{article}
\usepackage{graphicx}
\usepackage{algorithm}
\usepackage{algpseudocode}

\title{Formulazione sistema multi-agente}
\author{Lorenzo Pecorari}
\date{November 2025}
\maketitle
\begin{document}

\section{Spazio degli stati}
Sia S l'insieme degli stati tale che:
\begin{itemize}
    \item $S = \{s_0, s_1, ..., s_T\} = {\{s_t \}_{t \in [0, T]}, T \in \mathbb{N}_{>0}$
    \item $s_t = (s_t^1, s_t^2, ..., s_t^i, ..., s_t^N) = (s_t^i)_{i \in I}, I = \{1, 2, ..., N\}$
    \item $s_t^i = (b_t^i, \tau_t^i)$, stato del nodo $i$ al tempo $t$ 
    \item $b_t^i \in [0.0, 1.0]$ valore normalizzato di batteria residua di $i$ in $t$
    \item $\tau_t^i \in [0.0, 1.0]$, valore normalizzati di timestep completati da $i$ al tempo $t$ rispetto a $T$
\end{itemize}

\noindent Lo stato iniziale congiunto è definibile come $s_0 = (s_0^1, s_0^2, ..., s_0^N)$, dove ogni $s_0^i = (def\_batt, 0.0)$ con $def\_batt$ un valore definito dall'environment e attribuito ad ogni reset (inizio episodio) alla batteria.

\section{Spazio delle azioni}
Sia A l'insieme dello spazio delle azioni congiunte del sistema:
\begin{itemize}
    \item $A = \{a_1, a_2, ..., a_i, ..., a_K\}$
    \item $a_i = (a_i^1, a_i^2, ..., a_i^j, ..., a_i^N)$
    \item $a_i^j = (f_i^j, x_i^j, g_i^j, h_i^j)$ dove:
    \begin{itemize}
        \item $f_i^j \in [0, fps_{max}^j]$, framerate da adottare localmente per il nodo $j$ con valore massimo $fps_{max}^j$ definito a priori\\
        \item $x_i^j = $\{ \begin{array}{lr}
            $0 , ricezione o invio non abilitati$ \\
            $1 , invio lavoro ad altro nodo$\\
            $2 , ricezione lavoro da altro nodo$
            \end{array}\\
        \item $g_i^j = $\{ \begin{array}{lr}
            $0$, $\space se \space x_i^j$ = 0 \\
            $k \neq j$, se $x_i^j$ = 1 $\\
            \end{array}\\
        \item $h_i^j \in [0, fps_{max}^j]$, framerate per computazione di offloading
    \end{itemize}
\newpage
\end{itemize}

\newpage
\section{Nodo - agente}
\subsection{Parametri del nodo}
Definizione di parametri dei nodi:
\begin{itemize}
    \item $B$, capacità massima della batteria in Wh
    \item $P_I$, potenza assorbita in idle in W
    \item $P_M$, potenza assorbita a pieno carico in W
    \item $P_{PV}$, potenza in W generata dal pannello solare
    \item $A_{PV}$, superificie in $m^2$ del pannello
    \item $\eta_{PV}$, efficienza energetica del pannello
    \item $G(t)$, irradianza solare in $W/m^2$ al tempo t
    \item $G_{max}$, irradianza solare in $W/m^2$ massima
    \item $fps_{max}^i$, framerate massimo adaottabile dal nodo
    \item $\Delta t$, intervallo in s con cui simulare la computazione
    \item $E_{F, loc}$, energia in J consumata da per computare una frame localmente
    \item $E_{F, send}$, energia in J consumata per inviare frame ad un nodo per offloading
    \item $E_{F, recv}$, energia in J consumata per ricevere e computare frame da altro nodo
    \item $E_{B, t}$, energia in J della batteria al tempo t
    \item $E_{PV, t}$, energia in J proveniente dal pannello solare al tempo t
\end{itemize}

\subsubsection{Energia dal pannello}
Per ogni timestep $t \in [0, T]$, esiste un valore di $G(t) \in [0, G_{maz}]$ tale per cui si può definire $E_{PV, t} = \int{G(t)\cdot A_{PV} \cdot \eta_{PV} \space dt} \approx {G(t)\cdot A_{PV} \cdot \eta_{PV} \cdot \Delta t $, con $A_{PV} > 0$ e $\eta_{PV} \in[0.0, 1.0]$.

\subsubsection{Energia della batteria}
Per ogni timestep $t \in [0, T]$, la carica della batteria dipende dalla computazione effettuata in $t-1$. Ad ogni esecuzione di un'azione $a_t^i$, il relativo valore dell'energia residua nella batteria sarà determinato da\\\\
$E_{B, t}^i = (E_{B,t-1}^i + E_{PV,t}^i) - ((f_t^i \cdot E_{F, loc}^i \cdot \Delta t)+ (h_t^i \cdot E_{F, send/recv}^i \cdot \Delta t)) $

\subsubsection{Batteria e tempo normalizzati}
Per ogni timestep $t \in [0, T]$, la carica normalizzata della batteria $b_t^i$e i timestep completati $\tau_t^i$ normalizzati rispetto a T sono espressi da:
\begin{itemize}
    \item $b_t^i = \frac{E_{B,t}^i}{B^i}, B^i>0$
    \item $\tau_t^i = \frac{t}{T}, T > 0$
\end{itemize}


\subsection{Tabella Q}
Ogni nodo è caratterizzato da un agente che tiene conto di una propria Q table con $|S^i \times A^i|$ elementi. \\
Analogamente a quanto indicato in precedenza, ogni stato è definito come $s_t^i = (b_t^i, \tau_t^i)$, dove:
\begin{itemize}
    \item $b_t^i \in [0.0, 1.0]$, considerabili 11 valori differenti
    \item $\tau_t^i \in [0.0, 1.0]$, considerabili 11 valori differenti
\end{itemize}

Si assume quindi che $|S^i| = 11 \cdot 11 = 121$ .\\\\
In maniera simile allo spazio degli stati, è possibile definire lo spazio delle azioni caratterizzato dagli elementi $a_t^i = (f_t^i, x_t^i, g_t^i, h_t^i)$, dove:
\begin{itemize}
    \item $f_t^i \in [0, fps_{max}^i]$, considerabili $fps_{max}^i+1$ valori differenti
    \item $x_t^i \in \{0, 1, 2\}$, considerabili 3 valori differenti
    \item $g_t^i \in [0, n]$, considerabili $n+1$ valori con $n$ il numero di nodi
    \item $h_t^i \in [0, fps_{max}^i]$, considerabili $fps_{max}^i+1$ valori differenti
\end{itemize}

La dimensione dello spazio delle azioni è considerabile come:\\ $|A^i| = (fps_{max}^i+1) \cdot 3 \cdot (n+1) \cdot (fps_{max}^i+1) = 3(fps_{max}^i+1)^2\cdot(n+1)$\\\\
Date le due dimensioni degli spazi degli stati e delle azioni, la dimensione della Q table risulta essere: $|Q_{table}| = |S^i \times A^i| = |S^i| \cdot |A^i| = 121 \cdot 3(fps_{max}^i+1)^2\cdot(n+1) = 363(fps_{max}^i+1)^2(n+1)$\\

Supponendo che il sistema abbia $n = 2$ nodi ed entrambi abbiano $fps_{max} = 30$, la dimensione della tabella sarebbe di $1046529$ elementi. Al fine di ridurre tale valore, si potrebbe adottare una riduzione della dimensione di $f_t^i$ e $h_t^i$, in cui si considerano solo 5 possibili valori che permettono di variare di $fps_{max}^i / 5$ di volta in volta. Adottando tale assunzione, si ottiene che $|Q_{table}| = 363(fps_{max}^i+1)^2(n+1) \rightarrow 363(5)^2(2+1) = 81675$  elementi, una riduzione di 12.81 volte. 

\subsubsection{Scelta dell'azione}
Seguendo l'algoritmo $\epsilon-greedy$, si cerca di ottenere un trade-off tra exploration ed exploitation per l'agente, basandosi sul decay factor $\epsilon_{dec} \in[0, 1]$, sul valore iniziale $\epsilon_0$ e su quello finale $\epsilon_{fin}$. Per ogni timestep $t$, l'agente sceglie casualmente l'azione con probabilità $\epsilon$, mentre con probabilità $1 - \epsilon$ "effettua un ragionamento" in modo da ottimizzare la ricompensa.

\subsubsection{Osservazione dell'environment}
Ogni agente, ad ogni azione, riceve un reward dall'ambiente, il quale fornisce anche un'osservazione sullo stato. Al fine di non ottenere una tabella eccessivamente grande, si considera che ogni agente tenga conto solo del proprio stato al fine, quindi ciò che viene restituito dall'osservazione sarà solo $s_t^i = (b_t^i, \tau_t^i)$.

\subsubsection{Aggiornamento della tabella}
Ad ogni timestep $t$, l'agente sceglie un'azione $a_t^i$ sulla base dell'osservazione dello stato $s_t^i$ ottenuto dall'ambiente. Secondo la policy $\pi_{\theta}^i$, viene scelta un'azione $a_t^i = \pi_\theta^i(s_t^i)$, eseguita e collezionato il reward. La policy da seguire, in questo caso, prevede la scelta  dell'azione che massimizza il valore Q in corrispondenza di tale stato $s_t^i$: $a_t^i = argmax_{a \in A^i} Q(s_t^i, a)$.\\\\ A seguito di ciò si procede con l'aggiornamento della tabella secondo:\\

$Q(s_t^i, a_t^i) \leftarrow (1-\alpha)Q(s_t^i, a_t^i) + \alpha(r_t^i + \gamma(max_{a_{t+1}^i} Q(s_{t+1}^i, a_{t+1}^i)))$ \\\\
con $\alpha \in [0, 1]$ come learning rate e $\gamma \in [0, 1]$ come discount factor. 

\section{Funzione di reward}
Dati due nodi $i \in I$ e $j \in I$ tali che $j \neq i$, per ogni timestep $t \in [0, T]$, con azioni $a_t^i = (f_t^i, x_t^i, g_t^i, h_t^i)$ e $a_t^j = (f_t^j, x_t^j, g_t^j, h_t^j)$, è possibile definire la funzione di reward che permetta ad essi di ottenere una ricompensa sulla base dello stato corrente, dell'azione eseguita e dello stato successivo:\\

$R_i(t) = r_t^i = r_{frames, t}^i + r_{battery, t}^i + r_{cooperation, t}^i$ 
\begin{algorithm}
    \State{// Verifica se la somma degli fps adottati localmente e quelli di offloading sono minori uguali a quelli massimi per il nodo}
    \caption{Algoritmo per ricompensa per gestione delle frame, $r_{frames, t}^i$}\label{alg:cap}
    \begin{algoerithmic}
    \If{($(f_t^i + h_t^i) \leq fps_{max}^i$)}
        \State{return 1}
        %\State{return $\frac{(f_t^i + h_t^i)}{fps_{max}^i}$}
    \Else
    \State{return 0}
    \EndIf
    \end{algoerithmic}
\end{algorithm}

\begin{algorithm}
    \caption{Algoritmo per ricompensa per gestione della batteria, $r_{battery, t}^i$}\label{alg:cap}
    \begin{algoerithmic}
    \State{// Verifica se $i$ non effettua offloading}
    \If{($x_t^i =0$ \land $((f_t^i \cdot E_{F, loc} \cdot \Delta t) < (E_{B,t}^i + E_{PV, t}^i))$)}
        \State{return 1}\\
        %\State{return $(f_t^i \cdot E_{F, loc} \cdot \Delta t)$/ $(E_{B,t}^i + E_{PV, t}^i)$}\\

    \State{// Se $i$ invia lavoro, l'energia a disposizione deve essere necessaria}
    \ElsIf{($x_t^i =1$ \land $(((f_t^i \cdot E_{F, loc} \cdot \Delta t) + (h_t^i \cdot E_{F, send} \cdot \Delta t))) < (E_{B,t}^i + E_{PV, t}^i))$)}
        \State{return 1}\\
        
    %\State{return $((f_t^i \cdot E_{F, loc} \cdot \Delta t) + (h_t^i \cdot E_{F, send} \cdot \Delta t)) / (E_{B,t}^i + E_{PV, t}^i)$}\\


    \State{// Se $i$ riceve lavoro, l'energia a disposizione deve essere necessaria}
    \ElsIf{($x_t^i =2$ \land $(((f_t^i \cdot E_{F, loc} \cdot \Delta t) + (h_t^i \cdot E_{F, recv} \cdot \Delta t))) < (E_{B,t}^i + E_{PV, t}^i))$)}
        \State{return 1}\\
        %\State{return $((f_t^i \cdot E_{F, loc} \cdot \Delta t) + (h_t^i \cdot E_{F, recv} \cdot \Delta t)) / (E_{B,t}^i + E_{PV, t}^i)$}\\

    \State{// Vincolo non soddisfatto, ricompensa nulla}
    \Else
    \State{return 0}
    \EndIf
    \end{algoerithmic}
\end{algorithm}

\begin{algorithm}
    \caption{Algoritmo per ricompensa per cooperazione tra i nodi $r_{cooperation, t}^i$}\label{alg:cap}
    \begin{algoerithmic}

    \State{// Se $i$ non effettua offloading e non designa nessuno per tale operazione}
    \If{($x_t^i =0$ \land $g_t^i = 0$)}
        \State{return 1}\\

    \State{// Se $i$ effettua offloading, invia lavoro, designa un altro nodo e tale nodo intende ricevere da i}
    \ElsIf{($x_t^i =1$ \land $g_t^i \neq 0 \land g_t^i \neq i \land x_t^{g_t^i} = 2 \land g_t^{g_t^i} = i$)}
    \State{return 1}\\

    \State{// Se $i$ effettua offloading, riceve lavoro, designa un altro nodo e tale nodo intende inviare a i}
    \ElsIf{(($x_t^i =2$ \land $g_t^i \neq 0 \land g_t^i \neq i \land x_t^{g_t^i}= 1 \land g_t^{g_t^i} = i$))}
    \State{return 1}\\

    \State{// Nessun vincolo soddisfatto, ricompensa nulla}
    \Else
    \State{return 0}
    \EndIf
    \end{algoerithmic}
\end{algorithm}

\end{document}
